\documentclass[12pt,a4paper]{article}
\usepackage[utf8]{inputenc}
\usepackage[T5]{fontenc}
\usepackage[vietnamese]{babel}
\usepackage{geometry}
\usepackage{booktabs}
\usepackage{longtable}
\usepackage{enumitem}
\usepackage{titlesec}
\usepackage{hyperref}
\usepackage{xcolor}
\usepackage{fancyhdr}

\geometry{left=2cm,right=2cm,top=2cm,bottom=2cm}

% ============= ĐỊNH DẠNG TIÊU ĐỀ =============
\titleformat{\section}{\Large\bfseries\color{blue!70!black}}{}{0em}{\thesection.\ }
\titlespacing{\section}{0pt}{12pt plus 2pt minus 1pt}{6pt plus 1pt minus 1pt}

\titleformat{\subsection}{\large\bfseries\color{blue!60!black}}{}{0em}{\thesubsection.\ }
\titlespacing{\subsection}{0pt}{10pt plus 1pt minus 1pt}{4pt plus 0.5pt minus 0.5pt}

% ============= HEADER & FOOTER =============
\pagestyle{fancy}
\fancyhf{}
\fancyhead[L]{\textbf{Phân công nhiệm vụ -- Thiết kế Cơ sở dữ liệu}}
\fancyhead[R]{Nhóm 5 người -- 7 ngày}
\fancyfoot[C]{\thepage}
\renewcommand{\headrulewidth}{0.5pt}

% ============= TIÊU ĐỀ CHÍNH =============
\title{\textbf{\Large PHÂN CÔNG NHIỆM VỤ}\\[0.2em]\textbf{\Large THIẾT KẾ CƠ SỞ DỮ LIỆU}\\[0.5em]
{\normalsize Nền tảng kết nối cộng đồng nhiếp ảnh phim\\và các Film Lab tích hợp AI}\\[0.3em]
{\small\textit{(AI-powered Platform Connecting Film Photography Enthusiasts with Film Processing Labs)}}}
\author{Nhóm 5 thành viên}
\date{}

\begin{document}

\maketitle
\thispagestyle{empty}
\newpage

% =========================================================================
\section{Mục tiêu}
% =========================================================================

Thiết kế đầy đủ cơ sở dữ liệu PostgreSQL cho hệ thống, bao gồm:

\begin{itemize}[leftmargin=*]
    \item \textbf{Conceptual Design} (ERD)
    \item \textbf{Logical Design} (quan hệ, chuẩn hóa 3NF)
    \item \textbf{Physical Design} (kiểu dữ liệu, constraint, index)
    \item \textbf{Data Dictionary}
    \item \textbf{File SQL} hoàn chỉnh
\end{itemize}

Đảm bảo các thành viên \textbf{xen kẽ} công việc để hiểu rõ phần của nhau trong thời gian \textbf{7 ngày}.

\vspace{10pt}

% =========================================================================
\section{Phân công vai trò}
% =========================================================================

\begin{table}[h]
\centering
\begin{tabular}{@{}clp{5.5cm}p{5cm}@{}}
\toprule
\textbf{Thành viên} & \textbf{Vai trò} & \textbf{Phạm vi Ownership} & \textbf{Phải Review \& hiểu sâu} \\
\midrule
P1 & \textbf{Lead DB + Core Domain} & User -- Role -- FilmLab -- Order (trung tâm) & Toàn bộ quan hệ Order $\leftrightarrow$ Processing $\leftrightarrow$ Archive \\
P2 & \textbf{User \& Lab Management} & Auth, Profile, FilmLab, Service, Review, Rating & Order + Marketplace \\
P3 & \textbf{Order \& Processing} & Đặt hàng + Vòng đời xử lý phim & Digital Archive + Status History \\
P4 & \textbf{Digital Archive + Marketplace} & Kho ảnh số + Mua bán thiết bị & Order + Knowledge \\
P5 & \textbf{Knowledge + AI + Integration} & Community, Knowledge, AI, Logistics, Notification & Recommendation + Core User \\
\bottomrule
\end{tabular}
\end{table}

\vspace{10pt}

% =========================================================================
\section{Chi tiết nhiệm vụ từng thành viên}
% =========================================================================

\subsection{P1 -- Lead DB + Core Domain}

\begin{itemize}[leftmargin=*]
    \item Phân tích yêu cầu và xác định entity trung tâm.
    \item Thiết kế ERD tổng thể (Conceptual Design).
    \item Làm chủ các bảng: \texttt{users}, \texttt{roles}, \texttt{user\_roles}, \texttt{permissions}, \texttt{film\_labs}, \texttt{lab\_services}, \texttt{service\_packages}, \texttt{orders}, \texttt{order\_items}.
    \item Chịu trách nhiệm chuẩn hóa 3NF toàn bộ schema.
    \item Định nghĩa Business Rules và Integrity Constraints quan trọng.
    \item Tạo Data Dictionary tổng hợp.
    \item Quản lý version ERD và file SQL cuối cùng.
    \item Review sâu phần của P3 và P4.
\end{itemize}

\vspace{8pt}

\subsection{P2 -- User \& Film Lab Management}

\begin{itemize}[leftmargin=*]
    \item Thiết kế các entity: \texttt{user\_profiles}, \texttt{addresses}, \texttt{user\_preferences}, \texttt{film\_labs} (chi tiết), \texttt{lab\_working\_hours}, \texttt{lab\_capacity}, \texttt{lab\_services}, \texttt{service\_options}, \texttt{reviews}, \texttt{ratings}.
    \item Thiết kế phân quyền RBAC đầy đủ.
    \item Xử lý multi-tenant (mỗi Film Lab là một tenant).
    \item Định nghĩa các constraint: unique, check, foreign key.
    \item Review bắt buộc phần Order (P3) và Marketplace (P4).
\end{itemize}

\vspace{8pt}

\subsection{P3 -- Order \& Processing Lifecycle}

\begin{itemize}[leftmargin=*]
    \item Thiết kế các entity: \texttt{orders}, \texttt{order\_items}, \texttt{film\_rolls}, \texttt{processing\_stages}, \texttt{order\_status\_history}, \texttt{stage\_logs}, \texttt{pickup\_delivery\_requests}.
    \item Thiết kế state machine cho trạng thái đơn hàng.
    \item Xây dựng quan hệ chặt chẽ với Digital Archive (khi scan hoàn thành).
    \item Thiết kế index phục vụ truy vấn real-time tracking.
    \item Review bắt buộc phần Archive (P4) và Core (P1).
\end{itemize}

\vspace{8pt}

\subsection{P4 -- Digital Film Archive + Marketplace}

\begin{itemize}[leftmargin=*]
    \item \textbf{Digital Archive}: \texttt{scan\_files}, \texttt{photo\_metadata}, \texttt{albums}, \texttt{album\_photos}, \texttt{tags}, \texttt{photo\_tags}, \texttt{film\_stocks}, \texttt{camera\_models}, \texttt{lenses}.
    \item \textbf{Marketplace}: \texttt{products}, \texttt{product\_categories}, \texttt{listings}, \texttt{marketplace\_transactions}, \texttt{seller\_reviews}.
    \item Thiết kế metadata phong phú và cơ chế versioning file scan.
    \item Review bắt buộc phần Order (P3) và Knowledge (P5).
\end{itemize}

\vspace{8pt}

\subsection{P5 -- Knowledge + AI + Integration}

\begin{itemize}[leftmargin=*]
    \item \textbf{Knowledge \& Community}: \texttt{articles}, \texttt{tutorials}, \texttt{events}, \texttt{comments}, \texttt{knowledge\_categories}.
    \item \textbf{AI-related}: \texttt{recommendation\_logs}, \texttt{user\_behavior\_logs}, \texttt{knowledge\_chunks}, \texttt{embeddings} (metadata), \texttt{ai\_chat\_histories}, \texttt{scan\_quality\_scores}.
    \item \textbf{Integration}: \texttt{delivery\_partners}, \texttt{notifications}, \texttt{payment\_transactions}, \texttt{admin\_logs}.
    \item Thiết kế hỗ trợ semantic search và recommendation.
    \item Review bắt buộc phần User (P2) và Core (P1).
\end{itemize}

\vspace{10pt}

% =========================================================================
\section{Lịch trình 7 ngày (bắt buộc xen kẽ)}
% =========================================================================

\begin{description}[style=nextline, leftmargin=1.5cm, font=\bfseries\color{blue!60!black}]

    \item[\textbf{Ngày 1 -- Phân tích \& Conceptual Design (cả nhóm)}]
    \vspace{4pt}
    Đọc kỹ Functional Requirements và Core Flows; liệt kê toàn bộ Entity + Relationship; vẽ ERD sơ bộ chung; thống nhất naming convention, audit fields, soft-delete; chốt phân công chính thức.

    \vspace{8pt}

    \item[\textbf{Ngày 2 -- Logical Design phần riêng}]
    \vspace{4pt}
    Mỗi thành viên thiết kế đầy đủ Attribute, Domain, Primary Key, Foreign Key, Relationship và Constraint cho phần mình phụ trách.

    \vspace{8pt}

    \item[\textbf{Ngày 3 -- Cross-review lần 1 (bắt buộc)}]
    \vspace{4pt}
    \begin{itemize}[leftmargin=1.5em, font=\normalfont]
        \item P1 review P3 và P5
        \item P2 review P1 và P4
        \item P3 review P2 và P4
        \item P4 review P3 và P5
        \item P5 review P1 và P2
    \end{itemize}
    Mỗi người phải viết nhận xét và đề xuất thay đổi trực tiếp trên ERD/schema.

    \vspace{8pt}

    \item[\textbf{Ngày 4 -- Chỉnh sửa \& Quan hệ cross-module}]
    \vspace{4pt}
    Tập trung thiết kế các bảng trung gian và quan hệ phức tạp (Order $\leftrightarrow$ Archive, Lab $\leftrightarrow$ Review, User $\leftrightarrow$ AI\ldots).

    \vspace{8pt}

    \item[\textbf{Ngày 5 -- Normalization + Physical Design}]
    \vspace{4pt}
    Kiểm tra chuẩn hóa 3NF/BCNF; chọn kiểu dữ liệu PostgreSQL; thiết kế Index (đặc biệt search Lab theo vị trí, tracking order, full-text Knowledge); P5 dẫn dắt phần hỗ trợ AI.

    \vspace{8pt}

    \item[\textbf{Ngày 6 -- Gộp schema + Data Dictionary + Review lần 2}]
    \vspace{4pt}
    Tạo file SQL hoàn chỉnh (CREATE TABLE + CONSTRAINT + INDEX); viết Data Dictionary; thực hiện review xoay vòng lần 2; viết mẫu câu truy vấn quan trọng theo các Core Flow.

    \vspace{8pt}

    \item[\textbf{Ngày 7 -- Hoàn thiện \& Nộp bài}]
    \vspace{4pt}
    Hoàn thiện ERD cuối cùng (có màu theo module); file SQL PostgreSQL; Data Dictionary; Document thiết kế (giải thích các quyết định quan trọng); Checklist đối chiếu với Functional Requirements; P1 tổng hợp, cả nhóm xác nhận.

\end{description}

\vspace{10pt}

% =========================================================================
\section{Quy định bắt buộc để hiểu phần của nhau}
% =========================================================================

\begin{enumerate}[leftmargin=*]
    \item Mỗi thành viên \textbf{phải review ít nhất 2 phần} của người khác và viết nhận xét cụ thể.
    \item Mọi thay đổi phải được commit lên GitHub kèm theo lý do rõ ràng.
    \item Sử dụng chung một file ERD trực tuyến (dbdiagram.io hoặc draw.io).
    \item Cuối ngày 3 và ngày 6: mỗi thành viên trình bày ngắn phần mình đã review (5--7 phút).
    \item Không được chỉ hoàn thành phần mình mà bỏ qua phần còn lại của hệ thống.
\end{enumerate}

\vspace{15pt}

% =========================================================================
\begin{center}
\large\textit{Tài liệu này dùng làm căn cứ phân công và theo dõi tiến độ nhóm}\\
\textit{trong môn Thiết kế Cơ sở dữ liệu}
\end{center}
% =========================================================================

\end{document}